\chapter*{Abstract}
\thispagestyle{plain}

Named Entity Recognition (NER) findet benannte Entitäten wie Personen,
Orte oder Organisationen in Texten.
Encoder ordnen dazu jedem Token ein Label zu, Large Language Models
(LLMs) erzeugen die Entitäten als strukturierten Text.
Diese Arbeit vergleicht beide Wege: DeBERTa-v3-base und -large mit
vollständigem Fine-Tuning sowie Qwen3.5-0.8B, -4B und -27B im Zero-Shot-
und QLoRA-Regime auf dem englischen Teil von MultiNERD mit
15 Entitätstypen und automatisch erzeugter Referenzannotation.
Die fünf trainierten Konfigurationen werden mit je drei Seeds
ausgewertet.
Erfasst werden strikte Metriken auf Entitätsebene, die Formtreue der
JSON-Ausgaben, Trainingsdauer, Inferenzlatenz und GPU-Spitzenspeicher
der Inferenz; hinzu kommt ein Transfertest auf CrossNER.

DeBERTa-v3-large erreicht einen Mikro-F$_1$ von 95,8 Punkten,
Qwen3.5-27B im QLoRA-Regime zunächst 94,5.
Doch das Limit von 256 Ausgabetokens schneidet lange Antworten ab;
abgeschnittene zählen als leere Vorhersagen.
Die selektive Nachinferenz mit 1\,024 Tokens hebt das LLM auf 95,6 und
verringert den Abstand von 1,34 auf 0,18 Punkte (gepaartes
95-\%-Bootstrap-Intervall: 0,02 bis 0,33).
Im Makro-F$_1$ wird gegenüber dem großen Encoder kein eindeutiger
Unterschied nachgewiesen.
Auf CrossNER mit den gemeinsamen Typen PER, ORG und LOC kehrt sich die
Rangfolge bei höherem Limit um: Das LLM erreicht 87,1 Punkte, der Encoder
85,0 (Intervall der gepoolten Differenz: 1,47 bis 2,80 Punkte).

Die Encoder benötigen auf MultiNERD 9 bis 17~ms je Satz und 0,7 bis
1,7~GiB Inferenzspeicher, die QLoRA-Läufe mit 256 Ausgabetokens 1,9 bis
5,5~s und bis zu 18,7~GiB.
Auf gleicher Hardware ergibt sich für Qwen3.5-0.8B gegenüber
DeBERTa-v3-base ein Latenzfaktor von etwa 210.
Die Zero-Shot-Konfigurationen erreichen 6,2 bis 69,4 F$_1$-Punkte.
Prompt und Decoding dieser Läufe wurden nach Beobachtungen auf dem
Testsplit gewählt.
Da sich die Regime zudem über das Training hinaus in Prompt, Decoding und
numerischer Präzision unterscheiden, ist kein isolierter
Fine-Tuning-Effekt bestimmbar.

Auf MultiNERD erzielt bei jedem Latenz- und Speicherbudget, das eine
getestete Konfiguration zulässt, ein Encoder den höchsten Mikro-F$_1$.
Auf CrossNER benötigt das LLM rund 13~s je Satz und 18,6~GiB
Inferenzspeicher; ob der Qualitätsgewinn das rechtfertigt, hängt von der
Anwendung ab.
Der Qualitätsvergleich hängt jedoch wesentlich von Ausgabelimit,
Datensatz und der Typabbildung auf PER, ORG und LOC ab.
Die Arbeit zeigt damit zweierlei: den Ressourcenunterschied der konkreten
Pipelines und die Bedeutung ausreichender Ausgabelimits für die Bewertung
generativer NER.
Eine allgemeine Überlegenheit einer Modellklasse wird aus den
untersuchten Konfigurationen nicht abgeleitet.
